% ============================================================
% ГОСТ 2.104-2006 — Рамка и штампы для текстовых документов
% Подключается через % ============================================================
% ГОСТ 2.104-2006 — Рамка и штампы для текстовых документов
% Подключается через % ============================================================
% ГОСТ 2.104-2006 — Рамка и штампы для текстовых документов
% Подключается через % ============================================================
% ГОСТ 2.104-2006 — Рамка и штампы для текстовых документов
% Подключается через \input{gost_frame} в преамбуле документа
% Компилятор: XeLaTeX
% ============================================================

% ============================================================
% Данные штампа (РЕДАКТИРУЙТЕ ЗДЕСЬ)
% ============================================================
\newcommand{\GostDocDesignation}{ПМиИ.26.xx.yy.ПЗ}           % (2) Обозначение документа
\newcommand{\GostDocName}{Разработка нейросетевого сервиса для детекции БПЛА и~птиц с~интеграцией в~системы защиты наземных объектов}  % (1) Наименование
\newcommand{\GostOrganization}{Московский политех, гр.~221-171} % (9) Организация
\newcommand{\GostDeveloper}{Полх Т. А}                              % Разраб. — фамилия
\newcommand{\GostChecker}{Лукъянов М. Н.}                            % Пров.
\newcommand{\GostNormControl}{}                                % Н. контр.
\newcommand{\GostApprover}{}                                   % Утв.
\newcommand{\GostLitA}{}                                      % Литера 1
\newcommand{\GostLitB}{}                                       % Литера 2
\newcommand{\GostLitC}{}                                       % Литера 3
\newcommand{\GostMass}{}                                       % Масса
\newcommand{\GostScale}{}                                      % Масштаб
\newcommand{\GostSheetTotal}{\pageref{LastPage}}               % Листов

% ============================================================
% РАМКА — внешний прямоугольник
% 20мм слева, 5мм сверху/справа/снизу от края страницы
% ============================================================
\newcommand{\DrawGostFrame}{%
  \begin{tikzpicture}[remember picture, overlay]
    \coordinate (FrameBL) at ([xshift=20mm, yshift=5mm] current page.south west);
    \coordinate (FrameTR) at ([xshift=-5mm, yshift=-5mm] current page.north east);
    \draw[line width=0.7mm] (FrameBL) rectangle (FrameTR);
  \end{tikzpicture}%
}

% ============================================================
% ФОРМА 2 — основная надпись для текстовых документов
% (первый или заглавный лист)
% Размеры: 185 × 40 мм (8 строк по 5 мм)
% Столбцы: 7|10|23|15|10 | 70 | 15|15|20
% По ГОСТ 2.104-2006, стр. 9 (Форма 2)
% ============================================================
\newcommand{\DrawFormTwo}{%
  \begin{tikzpicture}[remember picture, overlay]
    % T0 = левый нижний угол штампа
    \coordinate (TBL) at ([xshift=-5mm, yshift=5mm] current page.south east);
    \coordinate (T0) at ([xshift=-185mm] TBL);
    %
    % --- Столбцы (x от T0, мм) ---
    % 0  7  17  40  55  65      135  150  165  185
    %  7  10  23  15  10   70    15   15   20
    \def\cA{0}    % левый край
    \def\cB{7}    % Изм.(7мм)
    \def\cC{17}   % Лист(10мм)
    \def\cD{40}   % № докум.(23мм)
    \def\cE{55}   % Подп.(15мм)
    \def\cF{65}   % Дата(10мм) → начало обл.(1)
    \def\cG{135}  % конец обл.(1) → начало правого блока
    \def\cH{150}  % Лит.(15мм) → Масса
    \def\cI{165}  % Масса(15мм) → Масштаб
    \def\cJ{185}  % Масштаб(20мм) → правый край
    \def\cK{140}     \def\cKa{145} 
    %
    % Подразделение Лит. на 3 ячейки по 5мм
    \def\cL{140}  % 135+5
    \def\cM{145}  % 135+10
    %
    % --- Строки (y от T0, снизу вверх, мм) ---
    % 0  5  10  15  20  25  30  35  40
    \def\rA{0}    % низ
    \def\rB{5}    % Утв.
    \def\rC{10}   % Н. контр.
    \def\rD{15}   % Пров.
    \def\rE{20}   % Разраб.
    \def\rF{25}   % Изм/Лист/etc + Лит/Масса/Масш заголовки
    \def\rG{30}
    \def\rH{35}
    \def\rI{40}   % верх
    %
    % =========================================
    % ГОРИЗОНТАЛЬНЫЕ ЛИНИИ
    % =========================================
    % Основные (полная ширина 0→185)
    \foreach \yy in {\rA, \rF, \rI} {
      \draw[line width=0.3mm]
        ([xshift=\cA mm, yshift=\yy mm] T0) --
        ([xshift=\cJ mm, yshift=\yy mm] T0);
    }
    %
    % Левый блок (0→65): строки данных
    \foreach \yy in {\rB, \rC, \rD, \rE, \rG, \rH} {
      \draw[line width=0.2mm]
        ([xshift=\cA mm, yshift=\yy mm] T0) --
        ([xshift=\cF mm, yshift=\yy mm] T0);
    }
    %
    % Правый блок (135→185): разделитель заголовки/значения Лит-Масса-Масш
    \foreach \yy in {\rE, \rD} {
    	\draw[line width=0.2mm]
    	([xshift=\cG mm, yshift=\yy mm] T0) --
    	([xshift=\cJ mm, yshift=\yy mm] T0);
    }


    % =========================================
    % ВЕРТИКАЛЬНЫЕ ЛИНИИ
    % =========================================
    % Внешние (полная высота)
    \foreach \xx in {\cA, \cJ} {
      \draw[line width=0.3mm]
        ([xshift=\xx mm, yshift=\rA mm] T0) --
        ([xshift=\xx mm, yshift=\rI mm] T0);
    }
    %
    \foreach \xx in {\cB} {
    	\draw[line width=0.2mm]
    	([xshift=\xx mm, yshift=\rF mm] T0) --
    	([xshift=\xx mm, yshift=\rI mm] T0);
    }
    % Левый блок: внутренние (0→25, под заголовком)
    \foreach \xx in {\cC, \cD, \cE, \cF} {
      \draw[line width=0.2mm]
        ([xshift=\xx mm, yshift=\rA mm] T0) --
        ([xshift=\xx mm, yshift=\rI mm] T0);
    }
    %
    % Центральный разделитель обл.(1) / правый блок (полная высота)
    \draw[line width=0.2mm]
      ([xshift=\cG mm, yshift=\rA mm] T0) --
      ([xshift=\cG mm, yshift=\rF mm] T0);
    %
    % Правый блок: Лит/Масса/Масштаб разделители (y=15→25)
    \foreach \xx in {\cI, \cH, \cJ} {
      \draw[line width=0.2mm]
        ([xshift=\xx mm, yshift=\rD mm] T0) --
        ([xshift=\xx mm, yshift=\rF mm] T0);
    }
    %
    
     \foreach \xx in {\cK, \cKa} {
    	\draw[line width=0.2mm]
    	([xshift=\xx mm, yshift=\rD mm] T0) --
    	([xshift=\xx mm, yshift=\rE mm] T0);
    }
    
    % =========================================
    % ТЕКСТ
    % =========================================
    \tikzset{
      lbl/.style={font=\fontsize{7pt}{8pt}\selectfont, inner sep=0pt},
      lblc/.style={lbl, anchor=center},
      lbll/.style={lbl, anchor=west},
      val/.style={font=\fontsize{9pt}{11pt}\selectfont, inner sep=0pt, anchor=center},
    }
    %
    % --- Заголовки столбцов (y=20→25) ---
    \node[lblc] at ([xshift=3.5mm,  yshift=27.5mm] T0) {Изм.};
    \node[lblc] at ([xshift=12mm,   yshift=27.5mm] T0) {Лист};
    \node[lblc] at ([xshift=28.5mm, yshift=27.5mm] T0) {№ докум.};
    \node[lblc] at ([xshift=47.5mm, yshift=27.5mm] T0) {Подп.};
    \node[lblc] at ([xshift=60mm,   yshift=27.5mm] T0) {Дата};
    %
    % --- Заголовки Лит./Масса/Масштаб (y=20→25) ---
    \node[lblc] at ([xshift=142.5mm, yshift=22.5mm] T0) {Лит.};
    \node[lblc] at ([xshift=157.5mm, yshift=22.5mm] T0) {Масса};
    \node[lblc] at ([xshift=175mm,   yshift=22.5mm] T0) {Масштаб};
    %
    % --- Значения Лит. (3 ячейки по 5мм) (y=15→20) ---
    \node[val] at ([xshift=137.5mm, yshift=17.5mm] T0) {\GostLitA};
    \node[val] at ([xshift=142.5mm, yshift=17.5mm] T0) {\GostLitB};
    \node[val] at ([xshift=147.5mm, yshift=17.5mm] T0) {\GostLitC};
    % Масса / Масштаб значения
    \node[val] at ([xshift=157.5mm, yshift=17.5mm] T0) {\GostMass};
    \node[val] at ([xshift=175mm,   yshift=17.5mm] T0) {\GostScale};
    %
    % --- Область (9) — организация (y=0→10, x=135→185) ---
    \node[val, text width=48mm, align=center]
      at ([xshift=160mm, yshift=7mm] T0) {\GostOrganization};
    %
    % --- Роли (y=0→20, x=0→17 объединено) ---
    \node[lbll] at ([xshift=1mm, yshift=22.5mm] T0) {Разраб.};
    \node[lbll] at ([xshift=1mm, yshift=17.5mm] T0) {Пров.};
    \node[lbll] at ([xshift=1mm, yshift=7.5mm]  T0) {Н. контр.};
    \node[lbll] at ([xshift=1mm, yshift=2.5mm]  T0) {Утв.};
    %
    % --- ФИО (x=17→40, col 2 = 23мм) ---
    \node[val] at ([xshift=28.5mm, yshift=22.5mm] T0) {\GostDeveloper};
    \node[val] at ([xshift=28.5mm, yshift=17.5mm] T0) {\GostChecker};
    \node[val] at ([xshift=28.5mm, yshift=7.5mm]  T0) {\GostNormControl};
    \node[val] at ([xshift=28.5mm, yshift=2.5mm]  T0) {\GostApprover};
    %
    % --- Область (2) — обозначение документа (y=25→40, x=65→185) ---
    \node[font=\fontsize{12pt}{14pt}\selectfont, anchor=center]
      at ([xshift=125mm, yshift=32.5mm] T0) {\GostDocDesignation};
    %
    % --- Область (1) — наименование (y=0→25, x=65→135) ---
    \node[font=\fontsize{11pt}{13pt}\selectfont, anchor=center,
          text width=65mm, align=center]
      at ([xshift=100mm, yshift=12.5mm] T0) {\GostDocName};
    %
  \end{tikzpicture}%
}

% ============================================================
% ФОРМА 2а — последующие листы
% Размеры: 185 × 15 мм (3 строки по 5 мм)
% Столбцы: 7|10|23|15|10 | 110 | 10
% По ГОСТ 2.104-2006, стр. 10 (Форма 2а)
% ============================================================
\newcommand{\DrawFormTwoA}{%
  \begin{tikzpicture}[remember picture, overlay]
    \coordinate (TBL) at ([xshift=-5mm, yshift=5mm] current page.south east);
    \coordinate (T0) at ([xshift=-185mm] TBL);
    %
    % Столбцы: 7|10|23|15|10|110|10 = 185
    \def\cA{0}   \def\cB{7}   \def\cC{17}  \def\cD{40}
    \def\cE{55}  \def\cF{65}  \def\cG{175} \def\cH{185}
    %
    % Внешний прямоугольник 185×15
    \draw[line width=0.3mm]
      (T0) rectangle ([xshift=185mm, yshift=15mm] T0);
    %
    % Горизонтальная линия: отделяет нижнюю строку (5мм) с Изм/Лист/...
    \draw[line width=0.2mm]
    
      ([yshift=5mm] T0) -- ([xshift=\cF mm, yshift=5mm] T0);
    \draw[line width=0.2mm]
      
      ([yshift=10mm] T0) -- ([xshift=\cF mm, yshift=10mm] T0);
    %
    % Вертикальные разделители левого блока (0→15мм, полная высота)
    \foreach \xx in {\cB, \cC, \cD, \cE} {
      \draw[line width=0.2mm]
        ([xshift=\xx mm] T0) -- ([xshift=\xx mm, yshift=15mm] T0);
    }
    %
    % Вертикальные разделители (полная высота 0→15мм)
    \foreach \xx in {\cF, \cG} {
      \draw[line width=0.2mm]
        ([xshift=\xx mm] T0) -- ([xshift=\xx mm, yshift=15mm] T0);
    }
    %
    % Текст
    \tikzset{
      lbl/.style={font=\fontsize{7pt}{8pt}\selectfont, inner sep=0pt, anchor=center},
    }
    \node[lbl] at ([xshift=3.5mm,  yshift=2.5mm] T0) {Изм.};
    \node[lbl] at ([xshift=12mm,   yshift=2.5mm] T0) {Лист};
    \node[lbl] at ([xshift=28.5mm, yshift=2.5mm] T0) {№ докум.};
    \node[lbl] at ([xshift=47.5mm, yshift=2.5mm] T0) {Подп.};
    \node[lbl] at ([xshift=60mm,   yshift=2.5mm] T0) {Дата};
    %
    % Область (2) — обозначение (x=65→175, y=0→15)
    \node[font=\fontsize{10pt}{12pt}\selectfont, anchor=center]
      at ([xshift=120mm, yshift=7.5mm] T0) {\GostDocDesignation};
    %
    % Горизонтальная линия в блоке «Лист» (отделяет надпись от номера)
    \draw[line width=0.2mm]
      ([xshift=175mm, yshift=8mm] T0) -- ([xshift=185mm, yshift=8mm] T0);
    %
    % Лист (номер страницы, x=175→185)
    \node[lbl] at ([xshift=180mm, yshift=11.5mm] T0) {Лист};
    \node[font=\fontsize{10pt}{12pt}\selectfont, anchor=center]
      at ([xshift=180mm, yshift=4mm] T0) {\thepage};
    %
  \end{tikzpicture}%
}

% ============================================================
% Стили страниц
% ============================================================
\fancypagestyle{gostfirst}{%
  \fancyhf{}%
  \renewcommand{\headrulewidth}{0pt}%
  \renewcommand{\footrulewidth}{0pt}%
  \fancyhead[C]{\DrawGostFrame}%
  \fancyfoot[C]{\DrawFormTwo}%
}

\fancypagestyle{gostplain}{%
  \fancyhf{}%
  \renewcommand{\headrulewidth}{0pt}%
  \renewcommand{\footrulewidth}{0pt}%
  \fancyhead[C]{\DrawGostFrame}%
  \fancyfoot[C]{\DrawFormTwoA}%
}

% Последующие страницы — Форма 2а
\pagestyle{gostplain}
 в преамбуле документа
% Компилятор: XeLaTeX
% ============================================================

% ============================================================
% Данные штампа (РЕДАКТИРУЙТЕ ЗДЕСЬ)
% ============================================================
\newcommand{\GostDocDesignation}{ПМиИ.26.xx.yy.ПЗ}           % (2) Обозначение документа
\newcommand{\GostDocName}{Разработка нейросетевого сервиса для детекции БПЛА и~птиц с~интеграцией в~системы защиты наземных объектов}  % (1) Наименование
\newcommand{\GostOrganization}{Московский политех, гр.~221-171} % (9) Организация
\newcommand{\GostDeveloper}{Полх Т. А}                              % Разраб. — фамилия
\newcommand{\GostChecker}{Лукъянов М. Н.}                            % Пров.
\newcommand{\GostNormControl}{}                                % Н. контр.
\newcommand{\GostApprover}{}                                   % Утв.
\newcommand{\GostLitA}{}                                      % Литера 1
\newcommand{\GostLitB}{}                                       % Литера 2
\newcommand{\GostLitC}{}                                       % Литера 3
\newcommand{\GostMass}{}                                       % Масса
\newcommand{\GostScale}{}                                      % Масштаб
\newcommand{\GostSheetTotal}{\pageref{LastPage}}               % Листов

% ============================================================
% РАМКА — внешний прямоугольник
% 20мм слева, 5мм сверху/справа/снизу от края страницы
% ============================================================
\newcommand{\DrawGostFrame}{%
  \begin{tikzpicture}[remember picture, overlay]
    \coordinate (FrameBL) at ([xshift=20mm, yshift=5mm] current page.south west);
    \coordinate (FrameTR) at ([xshift=-5mm, yshift=-5mm] current page.north east);
    \draw[line width=0.7mm] (FrameBL) rectangle (FrameTR);
  \end{tikzpicture}%
}

% ============================================================
% ФОРМА 2 — основная надпись для текстовых документов
% (первый или заглавный лист)
% Размеры: 185 × 40 мм (8 строк по 5 мм)
% Столбцы: 7|10|23|15|10 | 70 | 15|15|20
% По ГОСТ 2.104-2006, стр. 9 (Форма 2)
% ============================================================
\newcommand{\DrawFormTwo}{%
  \begin{tikzpicture}[remember picture, overlay]
    % T0 = левый нижний угол штампа
    \coordinate (TBL) at ([xshift=-5mm, yshift=5mm] current page.south east);
    \coordinate (T0) at ([xshift=-185mm] TBL);
    %
    % --- Столбцы (x от T0, мм) ---
    % 0  7  17  40  55  65      135  150  165  185
    %  7  10  23  15  10   70    15   15   20
    \def\cA{0}    % левый край
    \def\cB{7}    % Изм.(7мм)
    \def\cC{17}   % Лист(10мм)
    \def\cD{40}   % № докум.(23мм)
    \def\cE{55}   % Подп.(15мм)
    \def\cF{65}   % Дата(10мм) → начало обл.(1)
    \def\cG{135}  % конец обл.(1) → начало правого блока
    \def\cH{150}  % Лит.(15мм) → Масса
    \def\cI{165}  % Масса(15мм) → Масштаб
    \def\cJ{185}  % Масштаб(20мм) → правый край
    \def\cK{140}     \def\cKa{145} 
    %
    % Подразделение Лит. на 3 ячейки по 5мм
    \def\cL{140}  % 135+5
    \def\cM{145}  % 135+10
    %
    % --- Строки (y от T0, снизу вверх, мм) ---
    % 0  5  10  15  20  25  30  35  40
    \def\rA{0}    % низ
    \def\rB{5}    % Утв.
    \def\rC{10}   % Н. контр.
    \def\rD{15}   % Пров.
    \def\rE{20}   % Разраб.
    \def\rF{25}   % Изм/Лист/etc + Лит/Масса/Масш заголовки
    \def\rG{30}
    \def\rH{35}
    \def\rI{40}   % верх
    %
    % =========================================
    % ГОРИЗОНТАЛЬНЫЕ ЛИНИИ
    % =========================================
    % Основные (полная ширина 0→185)
    \foreach \yy in {\rA, \rF, \rI} {
      \draw[line width=0.3mm]
        ([xshift=\cA mm, yshift=\yy mm] T0) --
        ([xshift=\cJ mm, yshift=\yy mm] T0);
    }
    %
    % Левый блок (0→65): строки данных
    \foreach \yy in {\rB, \rC, \rD, \rE, \rG, \rH} {
      \draw[line width=0.2mm]
        ([xshift=\cA mm, yshift=\yy mm] T0) --
        ([xshift=\cF mm, yshift=\yy mm] T0);
    }
    %
    % Правый блок (135→185): разделитель заголовки/значения Лит-Масса-Масш
    \foreach \yy in {\rE, \rD} {
    	\draw[line width=0.2mm]
    	([xshift=\cG mm, yshift=\yy mm] T0) --
    	([xshift=\cJ mm, yshift=\yy mm] T0);
    }


    % =========================================
    % ВЕРТИКАЛЬНЫЕ ЛИНИИ
    % =========================================
    % Внешние (полная высота)
    \foreach \xx in {\cA, \cJ} {
      \draw[line width=0.3mm]
        ([xshift=\xx mm, yshift=\rA mm] T0) --
        ([xshift=\xx mm, yshift=\rI mm] T0);
    }
    %
    \foreach \xx in {\cB} {
    	\draw[line width=0.2mm]
    	([xshift=\xx mm, yshift=\rF mm] T0) --
    	([xshift=\xx mm, yshift=\rI mm] T0);
    }
    % Левый блок: внутренние (0→25, под заголовком)
    \foreach \xx in {\cC, \cD, \cE, \cF} {
      \draw[line width=0.2mm]
        ([xshift=\xx mm, yshift=\rA mm] T0) --
        ([xshift=\xx mm, yshift=\rI mm] T0);
    }
    %
    % Центральный разделитель обл.(1) / правый блок (полная высота)
    \draw[line width=0.2mm]
      ([xshift=\cG mm, yshift=\rA mm] T0) --
      ([xshift=\cG mm, yshift=\rF mm] T0);
    %
    % Правый блок: Лит/Масса/Масштаб разделители (y=15→25)
    \foreach \xx in {\cI, \cH, \cJ} {
      \draw[line width=0.2mm]
        ([xshift=\xx mm, yshift=\rD mm] T0) --
        ([xshift=\xx mm, yshift=\rF mm] T0);
    }
    %
    
     \foreach \xx in {\cK, \cKa} {
    	\draw[line width=0.2mm]
    	([xshift=\xx mm, yshift=\rD mm] T0) --
    	([xshift=\xx mm, yshift=\rE mm] T0);
    }
    
    % =========================================
    % ТЕКСТ
    % =========================================
    \tikzset{
      lbl/.style={font=\fontsize{7pt}{8pt}\selectfont, inner sep=0pt},
      lblc/.style={lbl, anchor=center},
      lbll/.style={lbl, anchor=west},
      val/.style={font=\fontsize{9pt}{11pt}\selectfont, inner sep=0pt, anchor=center},
    }
    %
    % --- Заголовки столбцов (y=20→25) ---
    \node[lblc] at ([xshift=3.5mm,  yshift=27.5mm] T0) {Изм.};
    \node[lblc] at ([xshift=12mm,   yshift=27.5mm] T0) {Лист};
    \node[lblc] at ([xshift=28.5mm, yshift=27.5mm] T0) {№ докум.};
    \node[lblc] at ([xshift=47.5mm, yshift=27.5mm] T0) {Подп.};
    \node[lblc] at ([xshift=60mm,   yshift=27.5mm] T0) {Дата};
    %
    % --- Заголовки Лит./Масса/Масштаб (y=20→25) ---
    \node[lblc] at ([xshift=142.5mm, yshift=22.5mm] T0) {Лит.};
    \node[lblc] at ([xshift=157.5mm, yshift=22.5mm] T0) {Масса};
    \node[lblc] at ([xshift=175mm,   yshift=22.5mm] T0) {Масштаб};
    %
    % --- Значения Лит. (3 ячейки по 5мм) (y=15→20) ---
    \node[val] at ([xshift=137.5mm, yshift=17.5mm] T0) {\GostLitA};
    \node[val] at ([xshift=142.5mm, yshift=17.5mm] T0) {\GostLitB};
    \node[val] at ([xshift=147.5mm, yshift=17.5mm] T0) {\GostLitC};
    % Масса / Масштаб значения
    \node[val] at ([xshift=157.5mm, yshift=17.5mm] T0) {\GostMass};
    \node[val] at ([xshift=175mm,   yshift=17.5mm] T0) {\GostScale};
    %
    % --- Область (9) — организация (y=0→10, x=135→185) ---
    \node[val, text width=48mm, align=center]
      at ([xshift=160mm, yshift=7mm] T0) {\GostOrganization};
    %
    % --- Роли (y=0→20, x=0→17 объединено) ---
    \node[lbll] at ([xshift=1mm, yshift=22.5mm] T0) {Разраб.};
    \node[lbll] at ([xshift=1mm, yshift=17.5mm] T0) {Пров.};
    \node[lbll] at ([xshift=1mm, yshift=7.5mm]  T0) {Н. контр.};
    \node[lbll] at ([xshift=1mm, yshift=2.5mm]  T0) {Утв.};
    %
    % --- ФИО (x=17→40, col 2 = 23мм) ---
    \node[val] at ([xshift=28.5mm, yshift=22.5mm] T0) {\GostDeveloper};
    \node[val] at ([xshift=28.5mm, yshift=17.5mm] T0) {\GostChecker};
    \node[val] at ([xshift=28.5mm, yshift=7.5mm]  T0) {\GostNormControl};
    \node[val] at ([xshift=28.5mm, yshift=2.5mm]  T0) {\GostApprover};
    %
    % --- Область (2) — обозначение документа (y=25→40, x=65→185) ---
    \node[font=\fontsize{12pt}{14pt}\selectfont, anchor=center]
      at ([xshift=125mm, yshift=32.5mm] T0) {\GostDocDesignation};
    %
    % --- Область (1) — наименование (y=0→25, x=65→135) ---
    \node[font=\fontsize{11pt}{13pt}\selectfont, anchor=center,
          text width=65mm, align=center]
      at ([xshift=100mm, yshift=12.5mm] T0) {\GostDocName};
    %
  \end{tikzpicture}%
}

% ============================================================
% ФОРМА 2а — последующие листы
% Размеры: 185 × 15 мм (3 строки по 5 мм)
% Столбцы: 7|10|23|15|10 | 110 | 10
% По ГОСТ 2.104-2006, стр. 10 (Форма 2а)
% ============================================================
\newcommand{\DrawFormTwoA}{%
  \begin{tikzpicture}[remember picture, overlay]
    \coordinate (TBL) at ([xshift=-5mm, yshift=5mm] current page.south east);
    \coordinate (T0) at ([xshift=-185mm] TBL);
    %
    % Столбцы: 7|10|23|15|10|110|10 = 185
    \def\cA{0}   \def\cB{7}   \def\cC{17}  \def\cD{40}
    \def\cE{55}  \def\cF{65}  \def\cG{175} \def\cH{185}
    %
    % Внешний прямоугольник 185×15
    \draw[line width=0.3mm]
      (T0) rectangle ([xshift=185mm, yshift=15mm] T0);
    %
    % Горизонтальная линия: отделяет нижнюю строку (5мм) с Изм/Лист/...
    \draw[line width=0.2mm]
    
      ([yshift=5mm] T0) -- ([xshift=\cF mm, yshift=5mm] T0);
    \draw[line width=0.2mm]
      
      ([yshift=10mm] T0) -- ([xshift=\cF mm, yshift=10mm] T0);
    %
    % Вертикальные разделители левого блока (0→15мм, полная высота)
    \foreach \xx in {\cB, \cC, \cD, \cE} {
      \draw[line width=0.2mm]
        ([xshift=\xx mm] T0) -- ([xshift=\xx mm, yshift=15mm] T0);
    }
    %
    % Вертикальные разделители (полная высота 0→15мм)
    \foreach \xx in {\cF, \cG} {
      \draw[line width=0.2mm]
        ([xshift=\xx mm] T0) -- ([xshift=\xx mm, yshift=15mm] T0);
    }
    %
    % Текст
    \tikzset{
      lbl/.style={font=\fontsize{7pt}{8pt}\selectfont, inner sep=0pt, anchor=center},
    }
    \node[lbl] at ([xshift=3.5mm,  yshift=2.5mm] T0) {Изм.};
    \node[lbl] at ([xshift=12mm,   yshift=2.5mm] T0) {Лист};
    \node[lbl] at ([xshift=28.5mm, yshift=2.5mm] T0) {№ докум.};
    \node[lbl] at ([xshift=47.5mm, yshift=2.5mm] T0) {Подп.};
    \node[lbl] at ([xshift=60mm,   yshift=2.5mm] T0) {Дата};
    %
    % Область (2) — обозначение (x=65→175, y=0→15)
    \node[font=\fontsize{10pt}{12pt}\selectfont, anchor=center]
      at ([xshift=120mm, yshift=7.5mm] T0) {\GostDocDesignation};
    %
    % Горизонтальная линия в блоке «Лист» (отделяет надпись от номера)
    \draw[line width=0.2mm]
      ([xshift=175mm, yshift=8mm] T0) -- ([xshift=185mm, yshift=8mm] T0);
    %
    % Лист (номер страницы, x=175→185)
    \node[lbl] at ([xshift=180mm, yshift=11.5mm] T0) {Лист};
    \node[font=\fontsize{10pt}{12pt}\selectfont, anchor=center]
      at ([xshift=180mm, yshift=4mm] T0) {\thepage};
    %
  \end{tikzpicture}%
}

% ============================================================
% Стили страниц
% ============================================================
\fancypagestyle{gostfirst}{%
  \fancyhf{}%
  \renewcommand{\headrulewidth}{0pt}%
  \renewcommand{\footrulewidth}{0pt}%
  \fancyhead[C]{\DrawGostFrame}%
  \fancyfoot[C]{\DrawFormTwo}%
}

\fancypagestyle{gostplain}{%
  \fancyhf{}%
  \renewcommand{\headrulewidth}{0pt}%
  \renewcommand{\footrulewidth}{0pt}%
  \fancyhead[C]{\DrawGostFrame}%
  \fancyfoot[C]{\DrawFormTwoA}%
}

% Последующие страницы — Форма 2а
\pagestyle{gostplain}
 в преамбуле документа
% Компилятор: XeLaTeX
% ============================================================

% ============================================================
% Данные штампа (РЕДАКТИРУЙТЕ ЗДЕСЬ)
% ============================================================
\newcommand{\GostDocDesignation}{ПМиИ.26.xx.yy.ПЗ}           % (2) Обозначение документа
\newcommand{\GostDocName}{Разработка нейросетевого сервиса для детекции БПЛА и~птиц с~интеграцией в~системы защиты наземных объектов}  % (1) Наименование
\newcommand{\GostOrganization}{Московский политех, гр.~221-171} % (9) Организация
\newcommand{\GostDeveloper}{Полх Т. А}                              % Разраб. — фамилия
\newcommand{\GostChecker}{Лукъянов М. Н.}                            % Пров.
\newcommand{\GostNormControl}{}                                % Н. контр.
\newcommand{\GostApprover}{}                                   % Утв.
\newcommand{\GostLitA}{}                                      % Литера 1
\newcommand{\GostLitB}{}                                       % Литера 2
\newcommand{\GostLitC}{}                                       % Литера 3
\newcommand{\GostMass}{}                                       % Масса
\newcommand{\GostScale}{}                                      % Масштаб
\newcommand{\GostSheetTotal}{\pageref{LastPage}}               % Листов

% ============================================================
% РАМКА — внешний прямоугольник
% 20мм слева, 5мм сверху/справа/снизу от края страницы
% ============================================================
\newcommand{\DrawGostFrame}{%
  \begin{tikzpicture}[remember picture, overlay]
    \coordinate (FrameBL) at ([xshift=20mm, yshift=5mm] current page.south west);
    \coordinate (FrameTR) at ([xshift=-5mm, yshift=-5mm] current page.north east);
    \draw[line width=0.7mm] (FrameBL) rectangle (FrameTR);
  \end{tikzpicture}%
}

% ============================================================
% ФОРМА 2 — основная надпись для текстовых документов
% (первый или заглавный лист)
% Размеры: 185 × 40 мм (8 строк по 5 мм)
% Столбцы: 7|10|23|15|10 | 70 | 15|15|20
% По ГОСТ 2.104-2006, стр. 9 (Форма 2)
% ============================================================
\newcommand{\DrawFormTwo}{%
  \begin{tikzpicture}[remember picture, overlay]
    % T0 = левый нижний угол штампа
    \coordinate (TBL) at ([xshift=-5mm, yshift=5mm] current page.south east);
    \coordinate (T0) at ([xshift=-185mm] TBL);
    %
    % --- Столбцы (x от T0, мм) ---
    % 0  7  17  40  55  65      135  150  165  185
    %  7  10  23  15  10   70    15   15   20
    \def\cA{0}    % левый край
    \def\cB{7}    % Изм.(7мм)
    \def\cC{17}   % Лист(10мм)
    \def\cD{40}   % № докум.(23мм)
    \def\cE{55}   % Подп.(15мм)
    \def\cF{65}   % Дата(10мм) → начало обл.(1)
    \def\cG{135}  % конец обл.(1) → начало правого блока
    \def\cH{150}  % Лит.(15мм) → Масса
    \def\cI{165}  % Масса(15мм) → Масштаб
    \def\cJ{185}  % Масштаб(20мм) → правый край
    \def\cK{140}     \def\cKa{145} 
    %
    % Подразделение Лит. на 3 ячейки по 5мм
    \def\cL{140}  % 135+5
    \def\cM{145}  % 135+10
    %
    % --- Строки (y от T0, снизу вверх, мм) ---
    % 0  5  10  15  20  25  30  35  40
    \def\rA{0}    % низ
    \def\rB{5}    % Утв.
    \def\rC{10}   % Н. контр.
    \def\rD{15}   % Пров.
    \def\rE{20}   % Разраб.
    \def\rF{25}   % Изм/Лист/etc + Лит/Масса/Масш заголовки
    \def\rG{30}
    \def\rH{35}
    \def\rI{40}   % верх
    %
    % =========================================
    % ГОРИЗОНТАЛЬНЫЕ ЛИНИИ
    % =========================================
    % Основные (полная ширина 0→185)
    \foreach \yy in {\rA, \rF, \rI} {
      \draw[line width=0.3mm]
        ([xshift=\cA mm, yshift=\yy mm] T0) --
        ([xshift=\cJ mm, yshift=\yy mm] T0);
    }
    %
    % Левый блок (0→65): строки данных
    \foreach \yy in {\rB, \rC, \rD, \rE, \rG, \rH} {
      \draw[line width=0.2mm]
        ([xshift=\cA mm, yshift=\yy mm] T0) --
        ([xshift=\cF mm, yshift=\yy mm] T0);
    }
    %
    % Правый блок (135→185): разделитель заголовки/значения Лит-Масса-Масш
    \foreach \yy in {\rE, \rD} {
    	\draw[line width=0.2mm]
    	([xshift=\cG mm, yshift=\yy mm] T0) --
    	([xshift=\cJ mm, yshift=\yy mm] T0);
    }


    % =========================================
    % ВЕРТИКАЛЬНЫЕ ЛИНИИ
    % =========================================
    % Внешние (полная высота)
    \foreach \xx in {\cA, \cJ} {
      \draw[line width=0.3mm]
        ([xshift=\xx mm, yshift=\rA mm] T0) --
        ([xshift=\xx mm, yshift=\rI mm] T0);
    }
    %
    \foreach \xx in {\cB} {
    	\draw[line width=0.2mm]
    	([xshift=\xx mm, yshift=\rF mm] T0) --
    	([xshift=\xx mm, yshift=\rI mm] T0);
    }
    % Левый блок: внутренние (0→25, под заголовком)
    \foreach \xx in {\cC, \cD, \cE, \cF} {
      \draw[line width=0.2mm]
        ([xshift=\xx mm, yshift=\rA mm] T0) --
        ([xshift=\xx mm, yshift=\rI mm] T0);
    }
    %
    % Центральный разделитель обл.(1) / правый блок (полная высота)
    \draw[line width=0.2mm]
      ([xshift=\cG mm, yshift=\rA mm] T0) --
      ([xshift=\cG mm, yshift=\rF mm] T0);
    %
    % Правый блок: Лит/Масса/Масштаб разделители (y=15→25)
    \foreach \xx in {\cI, \cH, \cJ} {
      \draw[line width=0.2mm]
        ([xshift=\xx mm, yshift=\rD mm] T0) --
        ([xshift=\xx mm, yshift=\rF mm] T0);
    }
    %
    
     \foreach \xx in {\cK, \cKa} {
    	\draw[line width=0.2mm]
    	([xshift=\xx mm, yshift=\rD mm] T0) --
    	([xshift=\xx mm, yshift=\rE mm] T0);
    }
    
    % =========================================
    % ТЕКСТ
    % =========================================
    \tikzset{
      lbl/.style={font=\fontsize{7pt}{8pt}\selectfont, inner sep=0pt},
      lblc/.style={lbl, anchor=center},
      lbll/.style={lbl, anchor=west},
      val/.style={font=\fontsize{9pt}{11pt}\selectfont, inner sep=0pt, anchor=center},
    }
    %
    % --- Заголовки столбцов (y=20→25) ---
    \node[lblc] at ([xshift=3.5mm,  yshift=27.5mm] T0) {Изм.};
    \node[lblc] at ([xshift=12mm,   yshift=27.5mm] T0) {Лист};
    \node[lblc] at ([xshift=28.5mm, yshift=27.5mm] T0) {№ докум.};
    \node[lblc] at ([xshift=47.5mm, yshift=27.5mm] T0) {Подп.};
    \node[lblc] at ([xshift=60mm,   yshift=27.5mm] T0) {Дата};
    %
    % --- Заголовки Лит./Масса/Масштаб (y=20→25) ---
    \node[lblc] at ([xshift=142.5mm, yshift=22.5mm] T0) {Лит.};
    \node[lblc] at ([xshift=157.5mm, yshift=22.5mm] T0) {Масса};
    \node[lblc] at ([xshift=175mm,   yshift=22.5mm] T0) {Масштаб};
    %
    % --- Значения Лит. (3 ячейки по 5мм) (y=15→20) ---
    \node[val] at ([xshift=137.5mm, yshift=17.5mm] T0) {\GostLitA};
    \node[val] at ([xshift=142.5mm, yshift=17.5mm] T0) {\GostLitB};
    \node[val] at ([xshift=147.5mm, yshift=17.5mm] T0) {\GostLitC};
    % Масса / Масштаб значения
    \node[val] at ([xshift=157.5mm, yshift=17.5mm] T0) {\GostMass};
    \node[val] at ([xshift=175mm,   yshift=17.5mm] T0) {\GostScale};
    %
    % --- Область (9) — организация (y=0→10, x=135→185) ---
    \node[val, text width=48mm, align=center]
      at ([xshift=160mm, yshift=7mm] T0) {\GostOrganization};
    %
    % --- Роли (y=0→20, x=0→17 объединено) ---
    \node[lbll] at ([xshift=1mm, yshift=22.5mm] T0) {Разраб.};
    \node[lbll] at ([xshift=1mm, yshift=17.5mm] T0) {Пров.};
    \node[lbll] at ([xshift=1mm, yshift=7.5mm]  T0) {Н. контр.};
    \node[lbll] at ([xshift=1mm, yshift=2.5mm]  T0) {Утв.};
    %
    % --- ФИО (x=17→40, col 2 = 23мм) ---
    \node[val] at ([xshift=28.5mm, yshift=22.5mm] T0) {\GostDeveloper};
    \node[val] at ([xshift=28.5mm, yshift=17.5mm] T0) {\GostChecker};
    \node[val] at ([xshift=28.5mm, yshift=7.5mm]  T0) {\GostNormControl};
    \node[val] at ([xshift=28.5mm, yshift=2.5mm]  T0) {\GostApprover};
    %
    % --- Область (2) — обозначение документа (y=25→40, x=65→185) ---
    \node[font=\fontsize{12pt}{14pt}\selectfont, anchor=center]
      at ([xshift=125mm, yshift=32.5mm] T0) {\GostDocDesignation};
    %
    % --- Область (1) — наименование (y=0→25, x=65→135) ---
    \node[font=\fontsize{11pt}{13pt}\selectfont, anchor=center,
          text width=65mm, align=center]
      at ([xshift=100mm, yshift=12.5mm] T0) {\GostDocName};
    %
  \end{tikzpicture}%
}

% ============================================================
% ФОРМА 2а — последующие листы
% Размеры: 185 × 15 мм (3 строки по 5 мм)
% Столбцы: 7|10|23|15|10 | 110 | 10
% По ГОСТ 2.104-2006, стр. 10 (Форма 2а)
% ============================================================
\newcommand{\DrawFormTwoA}{%
  \begin{tikzpicture}[remember picture, overlay]
    \coordinate (TBL) at ([xshift=-5mm, yshift=5mm] current page.south east);
    \coordinate (T0) at ([xshift=-185mm] TBL);
    %
    % Столбцы: 7|10|23|15|10|110|10 = 185
    \def\cA{0}   \def\cB{7}   \def\cC{17}  \def\cD{40}
    \def\cE{55}  \def\cF{65}  \def\cG{175} \def\cH{185}
    %
    % Внешний прямоугольник 185×15
    \draw[line width=0.3mm]
      (T0) rectangle ([xshift=185mm, yshift=15mm] T0);
    %
    % Горизонтальная линия: отделяет нижнюю строку (5мм) с Изм/Лист/...
    \draw[line width=0.2mm]
    
      ([yshift=5mm] T0) -- ([xshift=\cF mm, yshift=5mm] T0);
    \draw[line width=0.2mm]
      
      ([yshift=10mm] T0) -- ([xshift=\cF mm, yshift=10mm] T0);
    %
    % Вертикальные разделители левого блока (0→15мм, полная высота)
    \foreach \xx in {\cB, \cC, \cD, \cE} {
      \draw[line width=0.2mm]
        ([xshift=\xx mm] T0) -- ([xshift=\xx mm, yshift=15mm] T0);
    }
    %
    % Вертикальные разделители (полная высота 0→15мм)
    \foreach \xx in {\cF, \cG} {
      \draw[line width=0.2mm]
        ([xshift=\xx mm] T0) -- ([xshift=\xx mm, yshift=15mm] T0);
    }
    %
    % Текст
    \tikzset{
      lbl/.style={font=\fontsize{7pt}{8pt}\selectfont, inner sep=0pt, anchor=center},
    }
    \node[lbl] at ([xshift=3.5mm,  yshift=2.5mm] T0) {Изм.};
    \node[lbl] at ([xshift=12mm,   yshift=2.5mm] T0) {Лист};
    \node[lbl] at ([xshift=28.5mm, yshift=2.5mm] T0) {№ докум.};
    \node[lbl] at ([xshift=47.5mm, yshift=2.5mm] T0) {Подп.};
    \node[lbl] at ([xshift=60mm,   yshift=2.5mm] T0) {Дата};
    %
    % Область (2) — обозначение (x=65→175, y=0→15)
    \node[font=\fontsize{10pt}{12pt}\selectfont, anchor=center]
      at ([xshift=120mm, yshift=7.5mm] T0) {\GostDocDesignation};
    %
    % Горизонтальная линия в блоке «Лист» (отделяет надпись от номера)
    \draw[line width=0.2mm]
      ([xshift=175mm, yshift=8mm] T0) -- ([xshift=185mm, yshift=8mm] T0);
    %
    % Лист (номер страницы, x=175→185)
    \node[lbl] at ([xshift=180mm, yshift=11.5mm] T0) {Лист};
    \node[font=\fontsize{10pt}{12pt}\selectfont, anchor=center]
      at ([xshift=180mm, yshift=4mm] T0) {\thepage};
    %
  \end{tikzpicture}%
}

% ============================================================
% Стили страниц
% ============================================================
\fancypagestyle{gostfirst}{%
  \fancyhf{}%
  \renewcommand{\headrulewidth}{0pt}%
  \renewcommand{\footrulewidth}{0pt}%
  \fancyhead[C]{\DrawGostFrame}%
  \fancyfoot[C]{\DrawFormTwo}%
}

\fancypagestyle{gostplain}{%
  \fancyhf{}%
  \renewcommand{\headrulewidth}{0pt}%
  \renewcommand{\footrulewidth}{0pt}%
  \fancyhead[C]{\DrawGostFrame}%
  \fancyfoot[C]{\DrawFormTwoA}%
}

% Последующие страницы — Форма 2а
\pagestyle{gostplain}
 в преамбуле документа
% Компилятор: XeLaTeX
% ============================================================

% ============================================================
% Данные штампа (РЕДАКТИРУЙТЕ ЗДЕСЬ)
% ============================================================
\newcommand{\GostDocDesignation}{ПМиИ.26.xx.yy.ПЗ}           % (2) Обозначение документа
\newcommand{\GostDocName}{Разработка нейросетевого сервиса для детекции БПЛА и~птиц с~интеграцией в~системы защиты наземных объектов}  % (1) Наименование
\newcommand{\GostOrganization}{Московский политех, гр.~221-171} % (9) Организация
\newcommand{\GostDeveloper}{Полх Т. А}                              % Разраб. — фамилия
\newcommand{\GostChecker}{Лукъянов М. Н.}                            % Пров.
\newcommand{\GostNormControl}{}                                % Н. контр.
\newcommand{\GostApprover}{}                                   % Утв.
\newcommand{\GostLitA}{}                                      % Литера 1
\newcommand{\GostLitB}{}                                       % Литера 2
\newcommand{\GostLitC}{}                                       % Литера 3
\newcommand{\GostMass}{}                                       % Масса
\newcommand{\GostScale}{}                                      % Масштаб
\newcommand{\GostSheetTotal}{\pageref{LastPage}}               % Листов

% ============================================================
% РАМКА — внешний прямоугольник
% 20мм слева, 5мм сверху/справа/снизу от края страницы
% ============================================================
\newcommand{\DrawGostFrame}{%
  \begin{tikzpicture}[remember picture, overlay]
    \coordinate (FrameBL) at ([xshift=20mm, yshift=5mm] current page.south west);
    \coordinate (FrameTR) at ([xshift=-5mm, yshift=-5mm] current page.north east);
    \draw[line width=0.7mm] (FrameBL) rectangle (FrameTR);
  \end{tikzpicture}%
}

% ============================================================
% ФОРМА 2 — основная надпись для текстовых документов
% (первый или заглавный лист)
% Размеры: 185 × 40 мм (8 строк по 5 мм)
% Столбцы: 7|10|23|15|10 | 70 | 15|15|20
% По ГОСТ 2.104-2006, стр. 9 (Форма 2)
% ============================================================
\newcommand{\DrawFormTwo}{%
  \begin{tikzpicture}[remember picture, overlay]
    % T0 = левый нижний угол штампа
    \coordinate (TBL) at ([xshift=-5mm, yshift=5mm] current page.south east);
    \coordinate (T0) at ([xshift=-185mm] TBL);
    %
    % --- Столбцы (x от T0, мм) ---
    % 0  7  17  40  55  65      135  150  165  185
    %  7  10  23  15  10   70    15   15   20
    \def\cA{0}    % левый край
    \def\cB{7}    % Изм.(7мм)
    \def\cC{17}   % Лист(10мм)
    \def\cD{40}   % № докум.(23мм)
    \def\cE{55}   % Подп.(15мм)
    \def\cF{65}   % Дата(10мм) → начало обл.(1)
    \def\cG{135}  % конец обл.(1) → начало правого блока
    \def\cH{150}  % Лит.(15мм) → Масса
    \def\cI{165}  % Масса(15мм) → Масштаб
    \def\cJ{185}  % Масштаб(20мм) → правый край
    \def\cK{140}     \def\cKa{145} 
    %
    % Подразделение Лит. на 3 ячейки по 5мм
    \def\cL{140}  % 135+5
    \def\cM{145}  % 135+10
    %
    % --- Строки (y от T0, снизу вверх, мм) ---
    % 0  5  10  15  20  25  30  35  40
    \def\rA{0}    % низ
    \def\rB{5}    % Утв.
    \def\rC{10}   % Н. контр.
    \def\rD{15}   % Пров.
    \def\rE{20}   % Разраб.
    \def\rF{25}   % Изм/Лист/etc + Лит/Масса/Масш заголовки
    \def\rG{30}
    \def\rH{35}
    \def\rI{40}   % верх
    %
    % =========================================
    % ГОРИЗОНТАЛЬНЫЕ ЛИНИИ
    % =========================================
    % Основные (полная ширина 0→185)
    \foreach \yy in {\rA, \rF, \rI} {
      \draw[line width=0.3mm]
        ([xshift=\cA mm, yshift=\yy mm] T0) --
        ([xshift=\cJ mm, yshift=\yy mm] T0);
    }
    %
    % Левый блок (0→65): строки данных
    \foreach \yy in {\rB, \rC, \rD, \rE, \rG, \rH} {
      \draw[line width=0.2mm]
        ([xshift=\cA mm, yshift=\yy mm] T0) --
        ([xshift=\cF mm, yshift=\yy mm] T0);
    }
    %
    % Правый блок (135→185): разделитель заголовки/значения Лит-Масса-Масш
    \foreach \yy in {\rE, \rD} {
    	\draw[line width=0.2mm]
    	([xshift=\cG mm, yshift=\yy mm] T0) --
    	([xshift=\cJ mm, yshift=\yy mm] T0);
    }


    % =========================================
    % ВЕРТИКАЛЬНЫЕ ЛИНИИ
    % =========================================
    % Внешние (полная высота)
    \foreach \xx in {\cA, \cJ} {
      \draw[line width=0.3mm]
        ([xshift=\xx mm, yshift=\rA mm] T0) --
        ([xshift=\xx mm, yshift=\rI mm] T0);
    }
    %
    \foreach \xx in {\cB} {
    	\draw[line width=0.2mm]
    	([xshift=\xx mm, yshift=\rF mm] T0) --
    	([xshift=\xx mm, yshift=\rI mm] T0);
    }
    % Левый блок: внутренние (0→25, под заголовком)
    \foreach \xx in {\cC, \cD, \cE, \cF} {
      \draw[line width=0.2mm]
        ([xshift=\xx mm, yshift=\rA mm] T0) --
        ([xshift=\xx mm, yshift=\rI mm] T0);
    }
    %
    % Центральный разделитель обл.(1) / правый блок (полная высота)
    \draw[line width=0.2mm]
      ([xshift=\cG mm, yshift=\rA mm] T0) --
      ([xshift=\cG mm, yshift=\rF mm] T0);
    %
    % Правый блок: Лит/Масса/Масштаб разделители (y=15→25)
    \foreach \xx in {\cI, \cH, \cJ} {
      \draw[line width=0.2mm]
        ([xshift=\xx mm, yshift=\rD mm] T0) --
        ([xshift=\xx mm, yshift=\rF mm] T0);
    }
    %
    
     \foreach \xx in {\cK, \cKa} {
    	\draw[line width=0.2mm]
    	([xshift=\xx mm, yshift=\rD mm] T0) --
    	([xshift=\xx mm, yshift=\rE mm] T0);
    }
    
    % =========================================
    % ТЕКСТ
    % =========================================
    \tikzset{
      lbl/.style={font=\fontsize{7pt}{8pt}\selectfont, inner sep=0pt},
      lblc/.style={lbl, anchor=center},
      lbll/.style={lbl, anchor=west},
      val/.style={font=\fontsize{9pt}{11pt}\selectfont, inner sep=0pt, anchor=center},
    }
    %
    % --- Заголовки столбцов (y=20→25) ---
    \node[lblc] at ([xshift=3.5mm,  yshift=27.5mm] T0) {Изм.};
    \node[lblc] at ([xshift=12mm,   yshift=27.5mm] T0) {Лист};
    \node[lblc] at ([xshift=28.5mm, yshift=27.5mm] T0) {№ докум.};
    \node[lblc] at ([xshift=47.5mm, yshift=27.5mm] T0) {Подп.};
    \node[lblc] at ([xshift=60mm,   yshift=27.5mm] T0) {Дата};
    %
    % --- Заголовки Лит./Масса/Масштаб (y=20→25) ---
    \node[lblc] at ([xshift=142.5mm, yshift=22.5mm] T0) {Лит.};
    \node[lblc] at ([xshift=157.5mm, yshift=22.5mm] T0) {Масса};
    \node[lblc] at ([xshift=175mm,   yshift=22.5mm] T0) {Масштаб};
    %
    % --- Значения Лит. (3 ячейки по 5мм) (y=15→20) ---
    \node[val] at ([xshift=137.5mm, yshift=17.5mm] T0) {\GostLitA};
    \node[val] at ([xshift=142.5mm, yshift=17.5mm] T0) {\GostLitB};
    \node[val] at ([xshift=147.5mm, yshift=17.5mm] T0) {\GostLitC};
    % Масса / Масштаб значения
    \node[val] at ([xshift=157.5mm, yshift=17.5mm] T0) {\GostMass};
    \node[val] at ([xshift=175mm,   yshift=17.5mm] T0) {\GostScale};
    %
    % --- Область (9) — организация (y=0→10, x=135→185) ---
    \node[val, text width=48mm, align=center]
      at ([xshift=160mm, yshift=7mm] T0) {\GostOrganization};
    %
    % --- Роли (y=0→20, x=0→17 объединено) ---
    \node[lbll] at ([xshift=1mm, yshift=22.5mm] T0) {Разраб.};
    \node[lbll] at ([xshift=1mm, yshift=17.5mm] T0) {Пров.};
    \node[lbll] at ([xshift=1mm, yshift=7.5mm]  T0) {Н. контр.};
    \node[lbll] at ([xshift=1mm, yshift=2.5mm]  T0) {Утв.};
    %
    % --- ФИО (x=17→40, col 2 = 23мм) ---
    \node[val] at ([xshift=28.5mm, yshift=22.5mm] T0) {\GostDeveloper};
    \node[val] at ([xshift=28.5mm, yshift=17.5mm] T0) {\GostChecker};
    \node[val] at ([xshift=28.5mm, yshift=7.5mm]  T0) {\GostNormControl};
    \node[val] at ([xshift=28.5mm, yshift=2.5mm]  T0) {\GostApprover};
    %
    % --- Область (2) — обозначение документа (y=25→40, x=65→185) ---
    \node[font=\fontsize{12pt}{14pt}\selectfont, anchor=center]
      at ([xshift=125mm, yshift=32.5mm] T0) {\GostDocDesignation};
    %
    % --- Область (1) — наименование (y=0→25, x=65→135) ---
    \node[font=\fontsize{11pt}{13pt}\selectfont, anchor=center,
          text width=65mm, align=center]
      at ([xshift=100mm, yshift=12.5mm] T0) {\GostDocName};
    %
  \end{tikzpicture}%
}

% ============================================================
% ФОРМА 2а — последующие листы
% Размеры: 185 × 15 мм (3 строки по 5 мм)
% Столбцы: 7|10|23|15|10 | 110 | 10
% По ГОСТ 2.104-2006, стр. 10 (Форма 2а)
% ============================================================
\newcommand{\DrawFormTwoA}{%
  \begin{tikzpicture}[remember picture, overlay]
    \coordinate (TBL) at ([xshift=-5mm, yshift=5mm] current page.south east);
    \coordinate (T0) at ([xshift=-185mm] TBL);
    %
    % Столбцы: 7|10|23|15|10|110|10 = 185
    \def\cA{0}   \def\cB{7}   \def\cC{17}  \def\cD{40}
    \def\cE{55}  \def\cF{65}  \def\cG{175} \def\cH{185}
    %
    % Внешний прямоугольник 185×15
    \draw[line width=0.3mm]
      (T0) rectangle ([xshift=185mm, yshift=15mm] T0);
    %
    % Горизонтальная линия: отделяет нижнюю строку (5мм) с Изм/Лист/...
    \draw[line width=0.2mm]
    
      ([yshift=5mm] T0) -- ([xshift=\cF mm, yshift=5mm] T0);
    \draw[line width=0.2mm]
      
      ([yshift=10mm] T0) -- ([xshift=\cF mm, yshift=10mm] T0);
    %
    % Вертикальные разделители левого блока (0→15мм, полная высота)
    \foreach \xx in {\cB, \cC, \cD, \cE} {
      \draw[line width=0.2mm]
        ([xshift=\xx mm] T0) -- ([xshift=\xx mm, yshift=15mm] T0);
    }
    %
    % Вертикальные разделители (полная высота 0→15мм)
    \foreach \xx in {\cF, \cG} {
      \draw[line width=0.2mm]
        ([xshift=\xx mm] T0) -- ([xshift=\xx mm, yshift=15mm] T0);
    }
    %
    % Текст
    \tikzset{
      lbl/.style={font=\fontsize{7pt}{8pt}\selectfont, inner sep=0pt, anchor=center},
    }
    \node[lbl] at ([xshift=3.5mm,  yshift=2.5mm] T0) {Изм.};
    \node[lbl] at ([xshift=12mm,   yshift=2.5mm] T0) {Лист};
    \node[lbl] at ([xshift=28.5mm, yshift=2.5mm] T0) {№ докум.};
    \node[lbl] at ([xshift=47.5mm, yshift=2.5mm] T0) {Подп.};
    \node[lbl] at ([xshift=60mm,   yshift=2.5mm] T0) {Дата};
    %
    % Область (2) — обозначение (x=65→175, y=0→15)
    \node[font=\fontsize{10pt}{12pt}\selectfont, anchor=center]
      at ([xshift=120mm, yshift=7.5mm] T0) {\GostDocDesignation};
    %
    % Горизонтальная линия в блоке «Лист» (отделяет надпись от номера)
    \draw[line width=0.2mm]
      ([xshift=175mm, yshift=8mm] T0) -- ([xshift=185mm, yshift=8mm] T0);
    %
    % Лист (номер страницы, x=175→185)
    \node[lbl] at ([xshift=180mm, yshift=11.5mm] T0) {Лист};
    \node[font=\fontsize{10pt}{12pt}\selectfont, anchor=center]
      at ([xshift=180mm, yshift=4mm] T0) {\thepage};
    %
  \end{tikzpicture}%
}

% ============================================================
% Стили страниц
% ============================================================
\fancypagestyle{gostfirst}{%
  \fancyhf{}%
  \renewcommand{\headrulewidth}{0pt}%
  \renewcommand{\footrulewidth}{0pt}%
  \fancyhead[C]{\DrawGostFrame}%
  \fancyfoot[C]{\DrawFormTwo}%
}

\fancypagestyle{gostplain}{%
  \fancyhf{}%
  \renewcommand{\headrulewidth}{0pt}%
  \renewcommand{\footrulewidth}{0pt}%
  \fancyhead[C]{\DrawGostFrame}%
  \fancyfoot[C]{\DrawFormTwoA}%
}

% Последующие страницы — Форма 2а
\pagestyle{gostplain}
